\documentclass[11pt,a4paper]{article}

\usepackage[utf8]{inputenc}
\usepackage[T1]{fontenc}
\usepackage{lmodern}
\usepackage[top=0.75in, bottom=0.75in, left=0.85in, right=0.85in]{geometry}
\usepackage{hyperref}
\usepackage{parskip}

\hypersetup{
    colorlinks=true,
    linkcolor=blue,
    urlcolor=blue
}

\pagestyle{empty}

\begin{document}

% Candidate Header
{\large \textbf{\underline{Prabhavit Gupta}}}\\
Via Zanetti, 10, 30175 Venice, Italy \ $\cdot$ \ +39 351 610 7807 \ $\cdot$ \ \href{mailto:prabhavitg@gmail.com}{prabhavitg@gmail.com}\\
LinkedIn: \href{https://www.linkedin.com/in/prabhavit/}{https://www.linkedin.com/in/prabhavit}

\vspace{10pt}
September 28, 2026

\vspace{8pt}
\textbf{To:}\\
Selection Committee / Competency Centre Leads\\
Digital FAO and Agro-informatics Division (CSI)\\
Food and Agriculture Organization of the United Nations (FAO)\\
Viale delle Terme di Caracalla, 00153 Rome, Italy

\vspace{10pt}
\textbf{Subject: Expression of Interest: Information Technology Developer (Software Engineer / Web Developer / Data \& AI Engineer) -- CSI Division}

\vspace{6pt}
Dear Hiring Manager and Colleagues at FAO,

I am writing to express my strong interest in the Information Technology Developer opportunity within the Digital FAO and Agro-informatics Division (CSI). Having closely followed FAO's digital advancements, I am inspired by how the Agro-informatics unit combines robust software engineering, agricultural data pipelines, and AI-driven analysis to combat poverty and food insecurity. My background in full-stack web development, resilient data ingestion pipelines, and production AI orchestration directly aligns with CSI's mission to deliver secure, scalable, and impactful digital solutions for global operations.

Over the past three years, I have successfully transformed abstract business requirements into production-ready software features across multiple European technology projects spanning different countries and sectors; from a sports management platform in Luxembourg to strategic consulting platforms in France, Spain, and Germany. Throughout these engagements, I have prioritized pragmatic implementation, data integrity, and measurable business value within agile, multidisciplinary teams.

During my tenure at FlowSpace (London-based, remote), I managed our application ecosystem and led its modernization to a robust Next.js and Laravel stack, directly contributing to scaling the platform from 2,000 to over 4,500 active users. As part of this work, I designed and deployed an automated listing and data-ingestion pipeline (\texttt{listing.flowspace.co}) that converts heterogeneous operator sources---unstructured PDFs, spreadsheets, email payloads, and web data---into standardized, schema-validated marketplace assets. I implemented automated web scrapers with magic-byte file classification, established automated test suites (278 tests), created a human-in-the-loop review interface, and engineered an autonomous agent on GitHub Actions to auto-publish high-confidence records while routing edge cases for human review.

My daily engineering workflow maps directly onto the core profiles outlined in your announcement. In web and mobile development, I create responsive, accessible front-ends using TypeScript, React, Next.js, and Tailwind CSS, interfacing cleanly with backend REST APIs and microservices. In data and AI engineering, I build automated ETL workflows, design relational schemas in PostgreSQL and MySQL, and orchestrate LLM workflows using prompt chaining and RAG architectures connected to foundational models (Anthropic Claude, OpenAI). Furthermore, I adhere strictly to software development best practices: version control via Git and GitHub, defensive and secure programming, CI/CD pipeline automation, and clear architectural documentation.

Operating extensively across international projects, I have engineered automated document-parsing pipelines using n8n for UltraStrategy, automated multi-step workflows for Get2Germany, and executed rigorous quality assurance and API integrations for Clubee. These experiences have shaped my ability to ramp up rapidly inside unfamiliar codebases, communicate transparently across cross-functional teams, and maintain disciplined knowledge-sharing habits.

As a graduate of Ca' Foscari University of Venice (an IAU/UNESCO recognized institution), I have lived in Venice for over five years, hold an Italian EU Long-Term Residence Permit requiring no visa sponsorship anywhere in Italy or the European Union, and possess full professional English proficiency (C2). I am available with approximately two weeks' notice and am fully prepared for short- or medium-term travel to FAO Headquarters in Rome or field project locations as operational requirements demand.

I would welcome the opportunity to discuss how my practical background in full-stack engineering, automated data pipelines, and AI systems can immediately contribute to the Digital FAO and Agro-informatics Division.

Thank you for your time, consideration, and dedication to FAO's vital global mandate.

Sincerely,\\
\textbf{Prabhavit Gupta}

\vfill
\noindent\textit{\footnotesize I hereby authorize the processing of my personal data contained in this document in compliance with the European General Data Protection Regulation (EU Regulation 2016/679) and Italian Legislative Decree 196/2003, as modified by Legislative Decree 101/2018, for the purposes of recruitment and selection.}

\end{document}